% ---- pg-544 （书页 532） ----

3.\quad 合并 $U$ 与 $V$, 将新的类记为 $(UV)$. 将矩阵 $\mathbf{D}$ 中的元素更新: (a)删除原来 $U$ 和 $V$ 所对应的行和列, (b)将给出类 $(UV)$ 与其余各类之间距离的行和列加进该矩阵.

4.\quad 将步骤 2 和 3 重复 $N-1$ 次(当算法结束时所有对象将在一个类中). 记录下被合并的各个类的一致性以及历次合并的水平(距离或相似性).\hfill(12-12)

隐藏在任 何聚类方法背后的思想, 都可以通过实例来表达. 我们对连接方法的输入和算法组成部分作一番简要讨论之后再介绍这些例子.

\mvsub{单连接法}

单连接算法的输入可以是对象对之间的距离或相似性. 各个组由个体通过合并最近邻近点而形成, 此处的最近邻近点可以指最小距离或最大相似性.

一开始我们必须找到 $\mathbf{D}=\{d_{ij}\}$ 中的最小距离, 并将相应的对象比如说 $U$ 和 $V$ 合并, 得到聚类 $(UV)$. 对一般算法(12-12)中的步骤 3, $(UV)$ 与任一其他聚类 $W$ 之间的距离为
\phantomsection\label{eq:12-13}
\begin{equation}
d_{(UV)W}=\min\{d_{UW},d_{VW}\}
\tag{12-13}
\end{equation}
这里 $d_{UW}$ 和 $d_{VW}$ 分别为聚类 $U$ 与 $W$ 及聚类 $V$ 与 $W$ 的最近邻近点之间的距离.

单连接聚类的结果可以用树图表示. 树的分枝代表聚类. 分枝在结点处汇合(合并), 而沿距离(或相似性)轴的各个结点位置指明合并发生时的水平(即距离或相似性的大小). 我们在下面几个例子中考虑某些特殊情形下的树图.

\noindent\textbf{例 12.3}\quad (用单连接法聚类)

为说明单连接算法, 我们考虑以下的五个对象之间的假想距离如下:
\[
\mathbf{D}=\{d_{ik}\}=
\begin{array}{r@{\;}c}
 & \begin{array}{ccccc}1&2&3&4&5\end{array}\\[2pt]
\begin{array}{r}1\\2\\3\\4\\5\end{array} &
\left[\begin{array}{ccccc}
0&&&&\\
9&0&&&\\
3&7&0&&\\
6&5&9&0&\\
11&10&\tikz{\node[draw,ellipse,inner sep=1pt]{2};}&8&0
\end{array}\right]
\end{array}
\]
将每个对象看成是一个类, 我们的聚类从合并两个最接近的项目开始. 由于
\[
\min_{i,k}(d_{ik})=d_{53}=2
\]
对象 5 和 3 合并, 形成聚类(35). 为确定下一个聚类水平, 我们需要得到聚类(35)与其余对象 1, 2, 4 之间的距离. 最近邻近点距离为
\[
\begin{aligned}
d_{(35)1}&=\min\{d_{31},d_{51}\}=\min\{3,11\}=3\\
d_{(35)2}&=\min\{d_{32},d_{52}\}=\min\{7,10\}=7\\
d_{(35)4}&=\min\{d_{34},d_{54}\}=\min\{9,8\}=8
\end{aligned}
\]
从 $\mathbf{D}$ 中将对应于对象 3 和 5 的行和列删除, 并将对应于(35)的行和列加上, 便得到新的距离阵
\[
\begin{array}{r@{\;}c}
 & \begin{array}{cccc}(35)&1&2&4\end{array}\\[2pt]
\begin{array}{r}(35)\\1\\2\\4\end{array} &
\left[\begin{array}{cccc}
0&&&\\
\tikz{\node[draw,ellipse,inner sep=1pt]{3};}&0&&\\
7&9&0&\\
8&6&5&0
\end{array}\right]
\end{array}
\]
在这个距离阵中, 各对聚类之间的最小距离为 $d_{(35)1}=3$, 于是我们将聚类(1)与聚类(35)合并, 得到下一个聚类(135). 算出
